\documentclass[11pt]{article}

\usepackage[utf8]{inputenc}
\usepackage[T1]{fontenc}
\usepackage{amsmath,amssymb,amsthm}
\usepackage{mathtools}
\usepackage{bm}
\usepackage[margin=1in]{geometry}
\usepackage{booktabs}
\usepackage{array}
\usepackage{enumitem}
\usepackage{hyperref}
\usepackage{xcolor}
\hypersetup{colorlinks=true, linkcolor=blue!50!black, citecolor=blue!50!black,
            urlcolor=blue!50!black}

\theoremstyle{plain}
\newtheorem{theorem}{Theorem}
\newtheorem{lemma}{Lemma}
\theoremstyle{definition}
\newtheorem{assumption}{Assumption}
\newtheorem{remark}{Remark}

\newcommand{\Vx}{V_x}
\newcommand{\Vy}{V_y}
\newcommand{\Maug}{M_{\mathrm{aug}}}
\newcommand{\Deff}{D^{\mathrm{eff}}}
\newcommand{\Dwr}{D^{\mathrm{wr}}}
\newcommand{\sgn}{\operatorname{sgn}}
\newcommand{\Kmax}{K_{\max}}
\newcommand{\Tr}[1]{\tanh\!\left(#1\right)}

\title{\bfseries Velocity-Tracking Adaptive Sliding-Mode Controller:\\
Stability, Disturbance Headroom, and Initial-Gain Sizing\\
in the Julia Mecanum Simulator}
\author{\normalsize Mecanum\_MultiComponentFriction\_SlipSpin\_Roller\_ASMC\_3 (Julia)\\
\normalsize Mecanum-PINN v12.5 design note --- companion to the position-tracking analysis}
\date{June 9, 2026}

\begin{document}
\maketitle

\begin{abstract}
This note analyses the \emph{velocity-tracking} controller \texttt{asmc\_torques\_vel}
as the inner loop of a cascade, fed by an outer world-frame position/heading loop, and
uses the analysis to size the warm-start initial gains $K_{i0}$ that the velocity gain
law leaks toward. A first draft of this note treated the inner loop in isolation and
implicitly assumed body-frame velocity tracking alone certifies the world-frame
trajectory --- it does not. As written, \texttt{asmc\_torques\_vel} chases body-frame
setpoints that are open-loop in time; if the yaw rate tracks low, the heading error
$\Delta\psi$ accumulates, the body axes point the wrong way in the world, and the
world velocity $R(\psi)(\Vx,\Vy)^\top$ drifts even with all three body surfaces pinned
at zero. \S\ref{sec:frame} makes this failure explicit and shows why the cascade
resolves it: the outer loop measures world position/heading, sees the drift, and
re-commands the inner setpoints. The inner-loop relative-degree-one UUB proof
(\S\S\ref{sec:plant}--\ref{sec:uub}) is unchanged --- the surface dynamics collapse to
$\dot s_i=-\tilde K_i\Tr{s_i/\varepsilon_i}+\tilde d_i$ directly --- but the cascade
adds two inner-loop disturbance terms: a setpoint-rate term and a heading-rotation
(``circle'') term $\propto\dot\psi\,k_p\,e_{\mathrm{pos}}$, both quantified in
\S\ref{sec:cascade}. The heading correction acts on the \emph{shaped} error
$e_\psi(d_\psi)$ of the position controller \texttt{asmc\_torques}, not the raw
residual; \S\ref{sec:impl} gives the exact edits to \texttt{asmc\_torques\_vel}.
For deployment and run-time-matched PINN data, \S\ref{sec:mixed} develops the preferred
\emph{mixed relative-degree} variant --- degree-1 velocity surfaces on $(x,y)$, a
degree-2 position-type surface $s_\psi=\dot e_\psi+\lambda_\psi e_\psi$ on yaw --- which
bounds the heading error inside a single ASMC (no outer loop, no separation caveat) and
flags a $\sim38\,\dot\psi\,$N\,m yaw viscous term the present velocity controller omits. Timescale separation is comfortable: the inner loop runs at
$25$--$170\,$rad/s, leaving a $5\times$-slower outer loop at $\le5\,$rad/s
(translation) and $\le34\,$rad/s (yaw). Three numerical facts carry over: the
equivalent control inverts the \emph{augmented} $\Maug$ (cross-coupling
${\sim}13.4\%$, down from ${\sim}17.7\%$); the smaller velocity caps
$\Kmax=(60,80,100)\,$N\,m make actuator saturation the binding (B1) constraint at the
simultaneous-cap worst case; and the sign-indefinite leak-to-$K_0$ term is benign,
pinning the relaxed gain at $0.95K_{i0}$. The sizing rule
$K_{i0}\gtrsim\Dwr_{i,\mathrm{op}}/\tanh1$ now includes the cascade terms in
$\Dwr_{i,\mathrm{op}}$. All numerical values are verified against the exact platform
parameters in the notebook.
\end{abstract}

\tableofcontents

\section{What changes relative to the position controller}
\label{sec:changes}

The companion note analysed \texttt{asmc\_torques}, in which the sliding surface is
the standard first-order surface $s_i=\dot e_i+\lambda_i(e_i)\,e_i$ built on the
\emph{position} error, with a dynamic Gaussian $\lambda_i(e_i)$ and an
equivalent control that cancels the $\dot\lambda$ feedforward. The velocity
controller \texttt{asmc\_torques\_vel} is structurally simpler, and the
simplifications all act in the direction of a \emph{cleaner} proof:

\begin{enumerate}[leftmargin=2em]
\item \textbf{Relative degree one; no $\lambda$.} The surface is the body-frame
velocity error directly,
\begin{equation}
s_x=\Vx-\Vx^{d}, \qquad s_y=\Vy-\Vy^{d}, \qquad s_\psi=\dot\psi-\omega^{d},
\label{eq:surface}
\end{equation}
read off lines \texttt{s\_x = Vx - Vx\_d}, etc. There is no $\dot e+\lambda e$
construction, no dynamic-$\lambda$ Gaussian, and no $\dot\lambda$ term inside the
equivalent control. The whole of \S12 of the position note (``why the dynamic
$\lambda$ does not break the argument'') is vacuous here, and the
$\lambda$-cancellation that produced the clean $\dot s$ in the position note is
unnecessary --- the cancellation is automatic.

\item \textbf{Equivalent control is pure acceleration feedforward.} The set-points are
the body-frame desired accelerations directly,
$A_{x,\mathrm{eq}}=A_x^{d}$, $A_{y,\mathrm{eq}}=A_y^{d}$,
$\alpha_{\mathrm{eq}}=\alpha^{d}$, from \texttt{get\_Axb\_des}, \texttt{get\_Ayb\_des},
\texttt{get\_alpha\_des\_v}. There are \emph{no} error-transport terms
($-\lambda\dot e$, $-\dot\lambda e$) and no desired-velocity Coriolis correction
applied to the surface --- the Coriolis, viscous, and centripetal wrench terms in
$M_{x,\mathrm{eq}},M_{y,\mathrm{eq}},M_{\psi,\mathrm{eq}}$ are otherwise identical
to the position controller.

\item \textbf{The equivalent control inverts the augmented mass matrix.} The
inverse-dynamics wrench uses $\tilde m=m_s+4J_1/R^2$ and
$I_{\psi,\mathrm{aug}}=I_s+4(l+h)^2 J_1/R^2$ (wheel rotational inertia reflected to
the body), i.e. it inverts
\begin{equation}
\Maug=\begin{pmatrix}
\tilde m & 0 & -m a_Y\\
0 & \tilde m & m a_X\\
-m a_Y & m a_X & I_{\psi,\mathrm{aug}}
\end{pmatrix},
\label{eq:Maug}
\end{equation}
not the bare platform matrix $M$ of the position note. This is the matrix that
appears in \texttt{M\_aug} and in the MIMO-DOB compensation
$\Delta M_c=-\operatorname{diag}(R,R,1)\,\Maug\,\hat\delta$.

\item \textbf{A third gain-law term.} The velocity gain law is
\begin{equation}
\dot K_i=\underbrace{\gamma_i\,\sigma_i\,b(K_i)}_{\text{growth}}
-\underbrace{c_i\!\left(\tfrac{K_i}{K_{\max,i}}\right)^{\!3}}_{\text{cubic pushback}}
-\underbrace{\sigma^{\!*}_i\,(K_i-0.95\,K_{i0})\,
e^{\,\kappa\left(1-s_i^2/(9\varepsilon_\psi^2)\right)}}_{\text{leak toward }K_{i0}},
\label{eq:gainlaw}
\end{equation}
with $\sigma_i\equiv s_i\Tr{s_i/\varepsilon_i}\ge0$, smooth gate
$b(K)=\tfrac12-\tfrac12\Tr{K-(K_{\max}-2)}$, and the velocity configuration
($\sigma^{\!*}$ written for the leak coefficient to avoid clash with $\sigma_i$)
\begin{align*}
&\gamma=(8,\,15,\,25),\quad c=(0.1,\,0.3,\,0.5),\quad
K_{\max}=(60,\,80,\,100)\ \text{N\,m},\\
&\sigma^{\!*}=(0.5,\,0.25,\,0.25),\quad \kappa=0.25,\quad
K_{0}=(5,\,5,\,20)\ \text{N\,m},\quad
\varepsilon=0.05,\ \varepsilon_\psi=0.08.
\end{align*}
The leak term is new relative to the position note (which had only growth $+$ cubic).
It is sign-indefinite in $K_i-0.95K_{i0}$; we carry it through $\dot V$ in full in
\S\ref{sec:lyap}. Its fixed point is exactly the warm-start anchor $0.95\,K_{i0}$,
so sizing $K_{i0}$ is the practical content of this note.
\end{enumerate}

\section{What body-frame velocity tracking does \emph{not} certify}
\label{sec:frame}

The UUB result proven below bounds the three body-frame surfaces $s_x,s_y,s_\psi$.
It is tempting --- and a first draft of this note did exactly this --- to read that as
``the platform follows the intended trajectory.'' It does not. Two distinct gaps
must be separated.

\paragraph{Gap 1 (benign): static pose offset.} If the body-frame velocity is tracked
with a bounded but nonzero heading offset $\Delta\psi$, the world velocity is
$\dot r^{\mathrm{world}}=R(\psi)(\Vx,\Vy)^\top=R(\psi_{\mathrm{int}}-\Delta\psi)(\Vx,\Vy)^\top$,
a \emph{rotation} of the intended world velocity. The realised path is rotated, and
possibly translated, relative to the plan. For PINN training this is irrelevant: the
network is fed body-frame states $(\Vx,\Vy,\dot\psi,\omega_i)$ and forces, never world
position, so a rotated/translated trajectory is the same training datum. A bounded
$\Delta\psi$ is acceptable by design.

\paragraph{Gap 2 (the real problem): \emph{growing} heading error changes the
trajectory class.} The damaging case is a $\Delta\psi$ that \emph{grows with time}. A
circle authored as a constant body velocity $(\Vx^d,\Vy^d)=(v,0)$ with
$\omega^d=v/\rho$ becomes, under accumulating heading error, a spiral --- a
qualitatively different trajectory with different excitation content. This is the
concern that matters, and it lives entirely in the yaw channel.

\subsection{The heading error is the time-integral of the yaw surface}
The heading is the open-loop integral of yaw rate, and the intended heading is the
integral of the commanded yaw rate, so
\begin{equation}
\Delta\psi(t)\equiv\psi(t)-\psi_{\mathrm{int}}(t)
=\int_0^t\big(\dot\psi-\omega^d\big)\,d\tau=\int_0^t s_\psi(\tau)\,d\tau.
\label{eq:dpsi_integral}
\end{equation}
This is the crux: \textbf{heading error is the integral of the yaw sliding variable.}
UUB bounds $|s_\psi|\le\bar s_\psi$, but says nothing about
$\int s_\psi\,d\tau$. The worst-case bound is
\begin{equation}
|\Delta\psi(t)|\le\int_0^t|s_\psi|\,d\tau\le\bar s_\psi\,t,
\label{eq:dpsi_linear}
\end{equation}
which grows linearly. Whether $\Delta\psi$ actually drifts depends on the
\emph{mean} of the steady-state yaw residual:
\begin{itemize}[leftmargin=1.5em]
\item \textbf{Zero-mean $s_\psi$} (symmetric chatter about zero): the integral stays
bounded; $\Delta\psi$ wanders but does not secularly grow. Acceptable.
\item \textbf{Biased $s_\psi$} ($\langle s_\psi\rangle=\delta\neq0$): then
$\Delta\psi(t)\approx\delta\,t$ --- secular drift, circle $\to$ spiral. Not acceptable.
\end{itemize}

\subsection{What biases $\tilde d_\psi$, and hence $s_\psi$}
Inside the boundary layer the yaw surface settles at
\begin{equation}
s_\psi^{\mathrm{ss}}\approx\frac{\varepsilon_\psi\,\Dwr_\psi}{K_\psi},
\label{eq:spsi_ss}
\end{equation}
(the $\beta_\psi$ cancels through $\Deff_\psi=\beta_\psi\Dwr_\psi$, exactly as in
\eqref{eq:ss}). A \emph{constant} component of the yaw disturbance $\Dwr_\psi$
therefore produces a \emph{constant} $s_\psi^{\mathrm{ss}}$, which integrates into
linear heading drift. The yaw disturbance is generically \emph{not} zero-mean:
\begin{itemize}[leftmargin=1.5em]
\item the multicomponent spin torque $M_{zi}\propto\chi^2$ does not average to zero
over a sustained turn (the roller spin has a fixed sign while turning);
\item roller-handoff transients in $\tilde\varphi_i$ are phase-locked to wheel angle,
not symmetric in time;
\item most importantly, the DOB-onset \emph{structural asymmetry} already documented
(translation compensation gain ${\sim}2.25\times$ versus yaw lever
$L_\psi\approx0.13$) is a standing, signed offset on precisely this channel --- the
canonical DC bias source.
\end{itemize}

\subsection{Quantifying the open-loop drift}
With the operating yaw gain $K_\psi^{\mathrm{eq}}\approx79.4\,$N\,m (cubic balance,
\S\ref{sec:headroom}), $\varepsilon_\psi=0.08$, a representative biased yaw
disturbance gives the following linear drift rates:
\begin{center}
\begin{tabular}{lccc}
\toprule
Biased $\Dwr_\psi$ (N\,m) & 0.5 & 1.0 & 2.0\\
\midrule
$s_\psi^{\mathrm{ss}}$ (mrad/s) & 0.50 & 1.01 & 2.02\\
Heading drift over $30\,$s (deg) & 0.87 & 1.73 & 3.47\\
Heading error per circle lap (deg, $v{=}0.3$, $\rho{=}2$, $T{=}41.9\,$s) & 1.21 & 2.42 & 4.84\\
\bottomrule
\end{tabular}
\end{center}
The per-lap figure is the telling one: the error \emph{accumulates each lap}, so a
nominally closed circle opens into a spiral at $2$--$5^\circ$ per revolution. This is
exactly the failure mode of concern, and it is invisible to the inner-loop UUB
certificate.

\subsection{The cascade closes the integral, using the shaped heading error}
\label{sec:cascade}
The fix is to feed the inner yaw-rate setpoint from an outer heading loop that
\emph{measures} $\psi$ (which the simulator state carries at \texttt{state[4]}). To stay
consistent with the position controller \texttt{asmc\_torques}, the correction acts on
the \emph{shaped} heading error used there, not on the raw residual. With
$d_\psi\equiv\psi-\psi_d$, define
\begin{equation}
e_\psi(d_\psi)=2\tan\!\tfrac{d_\psi}{2}\,\big(1+2(1-\cos d_\psi)\big),
\qquad
e_\psi'(d_\psi)=\sec^2\!\tfrac{d_\psi}{2}\,(3-2\cos d_\psi)+4\tan\!\tfrac{d_\psi}{2}\sin d_\psi,
\label{eq:epsi}
\end{equation}
exactly the \texttt{e\_psi} and \texttt{edash\_psi} of the position controller. The
outer law is
\begin{equation}
\omega^d=\dot\psi_d+k_{p,\psi}\,e_\psi(d_\psi),
\qquad(\text{so } \omega^d>\dot\psi_d \text{ when } \psi \text{ leads } \psi_d).
\label{eq:outer}
\end{equation}
Three properties of \eqref{eq:epsi}, all verified numerically on $(-\pi,\pi)$, make this
a valid and well-behaved error coordinate:
\begin{enumerate}[leftmargin=2em]
\item \emph{Sign-faithful:} $\operatorname{sign}e_\psi(d_\psi)=\operatorname{sign}d_\psi$,
so the correction always pushes $\psi$ toward $\psi_d$.
\item \emph{Strictly monotone:} $e_\psi'(d_\psi)\ge e_\psi'(0)=1>0$ everywhere on the
domain, so there is no spurious equilibrium and the closed-loop pole never vanishes.
\item \emph{Identity at small angle:} $e_\psi(d_\psi)=d_\psi+O(d_\psi^3)$ with slope
exactly $1$ at the origin; the shaping only steepens at large excursions
($e_\psi/d_\psi\approx1.01$ at $0.1\,$rad, $\approx1.10$ at $0.3\,$rad), adding authority
where the heading is far off without altering the near-zero behaviour that sets the
steady-state drift.
\end{enumerate}
Differentiating \eqref{eq:dpsi_integral} (i.e. $\dot d_\psi=\dot\psi-\dot\psi_d$) and
substituting \eqref{eq:outer},
\begin{equation}
\dot d_\psi=\dot\psi-\omega^d
=\underbrace{(\dot\psi-\dot\psi_d)}_{\text{inner residual }s_\psi^{\mathrm{in}}}
-k_{p,\psi}\,e_\psi(d_\psi).
\label{eq:dpsi_closed}
\end{equation}
The accumulating integrator \eqref{eq:dpsi_integral} is replaced by a \emph{stable
nonlinear first-order} system. Linearising about the equilibrium $d_\psi^{\mathrm{ss}}$
gives the pole $-k_{p,\psi}\,e_\psi'(d_\psi^{\mathrm{ss}})$; since $e_\psi'\ge1$, the
shaped error is \emph{never slower} than a raw-error design and is up to ${\sim}30\%$
faster at large offsets. A biased inner residual no longer integrates without bound; it
settles where $k_{p,\psi}e_\psi(d_\psi^{\mathrm{ss}})=\langle s_\psi^{\mathrm{in}}\rangle$,
i.e.
\begin{equation}
e_\psi(d_\psi^{\mathrm{ss}})=\frac{\langle s_\psi^{\mathrm{in}}\rangle}{k_{p,\psi}}
=\frac{\varepsilon_\psi\,\Dwr_\psi}{k_{p,\psi}\,K_\psi}
\;\;\Longrightarrow\;\;
|d_\psi^{\mathrm{ss}}|\approx\frac{\varepsilon_\psi\,\Dwr_\psi}{k_{p,\psi}\,K_\psi}
\label{eq:dpsi_bound}
\end{equation}
where the final approximation uses $e_\psi\approx d_\psi$ at the small settled angle. The
resulting bounded heading offsets:
\begin{center}
\begin{tabular}{lcc}
\toprule
& $\Dwr_\psi=1.0$ & $\Dwr_\psi=2.0$\\
\midrule
$|d_\psi^{\mathrm{ss}}|$ at $k_{p,\psi}=2$ (deg) & 0.029 & 0.058\\
$|d_\psi^{\mathrm{ss}}|$ at $k_{p,\psi}=5$ (deg) & 0.012 & 0.023\\
$|d_\psi^{\mathrm{ss}}|$ at $k_{p,\psi}=10$ (deg) & 0.006 & 0.012\\
\bottomrule
\end{tabular}
\end{center}
The drift is now a fixed sub-tenth-of-a-degree offset (Gap 1, benign), not a secular
growth (Gap 2); the exact-$e_\psi$ form gives essentially the same steady bound as a raw
correction because $d_\psi^{\mathrm{ss}}$ is tiny, but inherits the position
controller's superior large-angle recovery. Optionally, the inner equivalent control can
cancel the rate of the correction, $\dot\omega^d=\ddot\psi_d+k_{p,\psi}e_\psi'(d_\psi)
\dot d_\psi$, by setting $\alpha_{\mathrm{eq}}=\alpha_{\mathrm{ff}}-k_{p,\psi}
e_\psi'(d_\psi)\,s_\psi$ (using $\dot d_\psi\approx s_\psi$); without it the un-cancelled
rate is the small setpoint-rate disturbance quantified below. The same construction
applies to the translational channels if world position must be bounded ---
$(\Vx^d,\Vy^d)\leftarrow R(\psi)^\top(\dot r_d^{\mathrm{world}}+K_p
e_{\mathrm{pos}}^{\mathrm{world}})$ --- but per the stated training use only the heading
channel needs closing; a bounded translational offset is acceptable.

\subsection{Timescale separation and the cost to the inner loop}
The cascade is valid only if the outer loop is slower than the inner. The inner
boundary-layer linearisation $\dot s_i\approx-(\beta_iK_i/\varepsilon_i)s_i$ gives
inner bandwidths
\begin{equation}
\omega^{\mathrm{in}}_i=\frac{\beta_iK_i^{\mathrm{eq}}}{\varepsilon_i}
\approx(24.8,\,28.2,\,171.3)\ \text{rad/s},
\end{equation}
so a $5\times$ separation permits outer gains up to
$\omega^{\mathrm{in}}_i/5\approx(5.0,\,5.6,\,34.3)\,$rad/s --- the yaw loop in
particular can be made fast. The cascade does add two terms to the inner-loop
disturbance, both small:
\begin{enumerate}[leftmargin=2em]
\item \emph{Setpoint-rate term.} The outer correction makes $\omega^d$ time-varying;
if the inner feedforward $\alpha_{\mathrm{eq}}$ does not include
$\dot\omega^d=k_{p,\psi}\,e_\psi'(d_\psi)\,\dot d_\psi$, the un-cancelled rate leaks into
$\tilde d_\psi$. Bounded by $k_{p,\psi}\,e_\psi'(d_\psi)\,\bar s_\psi$, i.e.\ a few
percent of the existing budget at the separations above (the $e_\psi'\to1$ near steady
state, so this is the same order as the raw-error estimate).
\item \emph{Heading-rotation (``circle'') term.} For the translational channels, if
closed, the rotation $R(\psi)^\top$ applied to a world correction injects a body-frame
setpoint rate $\dot\psi\,k_p\,e_{\mathrm{pos}}$; at $\dot\psi=0.3\,$rad/s,
$k_p=1$, $e_{\mathrm{pos}}=0.1\,$m this is $0.03\,$m/s$^2$, i.e.\ ${\sim}1.35\,$N\,m of
wrench --- negligible against the $\sim$45--76\,N\,m headroom.
\end{enumerate}
Both are folded into $\Dwr_{i,\mathrm{op}}$ for the $K_{i0}$ sizing of
\S\ref{sec:K0}; neither threatens the (B1) margin.

\subsection{Implementation in \texttt{asmc\_torques\_vel}}
\label{sec:impl}
The cascade requires three local edits to \texttt{asmc\_torques\_vel} and one new field
on \texttt{ASMCParams}; the inner relative-degree-one structure is untouched.
\begin{enumerate}[leftmargin=2em]
\item \textbf{Unpack the heading.} Add $\psi=\texttt{state[4]}$ to the state unpack
(currently only \texttt{state[1:3,17:19]} are read).
\item \textbf{Corrected yaw-rate setpoint.} Replace the open-loop $\omega_d$ with the
shaped-error cascade \eqref{eq:outer}: with $\omega_{\mathrm{ff}}=
\texttt{get\_omega\_des\_v(t)}$, $\psi_d=\texttt{get\_psi\_des(t)}$,
$d_\psi=\psi-\psi_d$, $e_\psi$ from \eqref{eq:epsi},
\[
\omega_d=\omega_{\mathrm{ff}}+\texttt{asmc.kp\_psi}\cdot e_\psi(d_\psi),
\qquad s_\psi=\dot\psi-\omega_d .
\]
\item \textbf{(Optional) feed-forward the correction rate.} Replace
$\alpha_{\mathrm{eq}}=\texttt{get\_alpha\_des\_v(t)}$ with
$\alpha_{\mathrm{eq}}=\alpha_{\mathrm{ff}}-\texttt{asmc.kp\_psi}\cdot
e_\psi'(d_\psi)\,s_\psi$, using $e_\psi'$ from \eqref{eq:epsi}. This cancels the
setpoint-rate disturbance; omitting it leaves the small leak of item~1 above.
\end{enumerate}
The new parameter is a single field \texttt{kp\_psi::Float64 = 5.0} (rad/s) on
\texttt{ASMCParams}. The choice $k_{p,\psi}=5$ sits at $\sim\!\omega^{\mathrm{in}}_\psi/34$,
far inside the $5\times$ separation margin, and yields $|d_\psi^{\mathrm{ss}}|\lesssim
0.02^\circ$. Two consistency requirements on the reference, both satisfied by the
existing symbolic construction ($\psi_{\mathrm{sym}}=\pi/2+\psi_{\mathrm{silu}}$,
$\omega_{\mathrm{sym}}=\dot\psi_{\mathrm{sym}}$): (i) \texttt{get\_psi\_des} must be the
exact integral of the yaw-rate feedforward that \texttt{get\_omega\_des\_v} returns, or
the correction fights the feedforward; (ii) at $t=0$, $\psi(0)=\psi_d(0)$ makes
$e_\psi(0)=0$, so the correction is inert at start and adds no startup transient.

\section{Plant on the sliding manifold and the canonical form}
\label{sec:plant}

With $v=(\Vx,\Vy,\dot\psi)^\top$, the body-frame dynamics actually used by the
velocity controller read
\begin{equation}
\Maug\,\dot v=\tau_{\mathrm{task}}+\mathcal{F}_{\mathrm{fric}}(\chi,\mu,v,\omega,\theta)
+\mathcal{C}(v),
\label{eq:plant}
\end{equation}
where $\tau_{\mathrm{task}}=(M_x^{\mathrm{cmd}},M_y^{\mathrm{cmd}},
M_\psi^{\mathrm{cmd}})^\top$ is the wrench delivered after the wheel-map and
$\tanh$ saturation, and $\mathcal{F}_{\mathrm{fric}}$ collects the multicomponent
LuGre/Adamov forces and spin torques. Because the equivalent control is precisely
$\Maug A_{\mathrm{eq}}-\mathcal{C}(v)-\mathcal{F}^{\mathrm{model}}_{\mathrm{fric}}$,
the closed loop on the nominal model is
\begin{equation}
\dot v=A_{\mathrm{eq}}+\Maug^{-1}\tau_{\mathrm{sw}}+\Maug^{-1}d(t),
\label{eq:closed}
\end{equation}
where $d(t)$ lumps friction-model mismatch (the $\chi$-dependent spin torque
$M_{zi}$ in particular), parametric error in $(m_s,I_s,J_1,p_1,p_2,a_X,a_Y)$,
actuator-saturation residual, roller-handoff transients, and any state-estimation
error if run open-loop on a PINN.

Differentiating the surface \eqref{eq:surface} and substituting
\eqref{eq:closed}, with $A_{\mathrm{eq}}=\dot v^{d}$ (the desired body acceleration),
\begin{equation}
\dot s_i=[\dot v]_i-[\dot v^{d}]_i
=[\Maug^{-1}\tau_{\mathrm{sw}}]_i+[\Maug^{-1}d]_i .
\label{eq:sdot}
\end{equation}
There is no pairwise $\lambda$-cancellation to perform: relative degree one means the
feedforward $A_{\mathrm{eq}}=\dot v^{d}$ cancels the desired-acceleration term
exactly, and \eqref{eq:sdot} is immediate. This is the structural payoff of velocity
tracking --- the same canonical $\dot s$ the position note had to \emph{engineer}
via $\dot\lambda$ inside $A_{i,\mathrm{eq}}$ falls out for free.

\begin{remark}[Channel coupling through $\Maug^{-1}$; the velocity numbers]
\label{rem:coupling}
With the exact notebook parameters ($h=0.235$, $l=0.15$, $R=0.05$, $R_d=0.0355$,
$m=30$, $m_w=1.4$, $J_1=5.87\times10^{-3}$, $m_s=35.6$, $I_s=4.42$, so
$\tilde m=44.99$, $I_{\psi,\mathrm{aug}}=5.812$, $m a_X=0.48$, $-m a_Y=0.78$):
\begin{equation}
\Maug=\begin{pmatrix}44.992&0&0.780\\0&44.992&0.480\\0.780&0.480&5.812\end{pmatrix},
\quad
\Maug^{-1}\!\approx\!\begin{pmatrix}
2.228\!\times\!10^{-2}&3.19\!\times\!10^{-5}&-2.992\!\times\!10^{-3}\\
3.19\!\times\!10^{-5}&2.225\!\times\!10^{-2}&-1.841\!\times\!10^{-3}\\
-2.992\!\times\!10^{-3}&-1.841\!\times\!10^{-3}&1.726\!\times\!10^{-1}
\end{pmatrix}.
\end{equation}
The smallest eigenvalue of $\Maug$ is $\approx5.79>0$, so $\Maug$ is symmetric
positive definite and the diagonal projection coefficients
$\beta_i\le[\Maug^{-1}]_{ii}$ are bounded away from zero. As fractions of the
diagonal,
\[
\frac{\sum_{j\neq x}|[\Maug^{-1}]_{xj}|}{[\Maug^{-1}]_{xx}}\approx13.6\%,\quad
\frac{\sum_{j\neq y}|[\Maug^{-1}]_{yj}|}{[\Maug^{-1}]_{yy}}\approx8.4\%,\quad
\frac{\sum_{j\neq\psi}|[\Maug^{-1}]_{\psi j}|}{[\Maug^{-1}]_{\psi\psi}}\approx2.8\%.
\]
The worst coupling is again $x\!\leftrightarrow\!\psi$, but at $13.4\%$ it is
\emph{lower} than the position note's $17.7\%$: reflecting the wheel inertia into the
body inflates the diagonal $\tilde m$ and $I_{\psi,\mathrm{aug}}$ faster than the
off-diagonals $m a_X,m a_Y$ (which are unchanged), shrinking the relative coupling.
The two devices of the position note carry over verbatim:
\begin{itemize}[leftmargin=1.5em]
\item \emph{Device 1.} Absorb cross-axis switching into the disturbance. Since
$|\Tr{s_j/\varepsilon_j}|\le1$ and the gate keeps $K_j\le\Kmax+O(c/\gamma)$,
\begin{equation}
\Deff_i=\Dwr_i\!\cdot\!\beta_i+\sum_{j\neq i}|[\Maug^{-1}]_{ij}|\,K_{\max,j},
\label{eq:Deff}
\end{equation}
where $\Dwr_i$ is the genuine model-mismatch in wrench units.
\item \emph{Device 2.} Define $\beta_x\le0.0223$, $\beta_y\le0.0222$,
$\beta_\psi\le0.1726$ (a safe practical choice $\beta_x=\beta_y=0.020$,
$\beta_\psi=0.15$). With $\tilde K_i=\beta_i K_i$,
\begin{equation}
\dot s_i=-\tilde K_i\Tr{s_i/\varepsilon_i}+\tilde d_i,\qquad
|\tilde d_i|\le\Deff_i,
\label{eq:canonical}
\end{equation}
the canonical scalar form.
\end{itemize}
\end{remark}

\section{Lyapunov function and its derivative}
\label{sec:lyap}

\subsection{Candidate}
Identical to the position note. For each axis fix a nominal ideal gain
$K_i^*>0$ and use the gain-error rescaling that places $\beta_i$ in the numerator,
\begin{equation}
V_i(s_i,K_i)=\tfrac12 s_i^2+\frac{\beta_i}{2\gamma_i}(K_i-K_i^*)^2,
\qquad
V=\sum_i V_i.
\label{eq:V}
\end{equation}

\subsection{Derivative, carrying the leak term in full}
Differentiating \eqref{eq:V} along \eqref{eq:canonical} and \eqref{eq:gainlaw}, with
$\sigma_i\equiv s_i\Tr{s_i/\varepsilon_i}\ge0$ and writing
$g_i\equiv e^{\,\kappa(1-s_i^2/(9\varepsilon_\psi^2))}>0$ for the (bounded, positive)
leak modulation:
\begin{align}
\dot V_i
&= s_i\dot s_i+\frac{\beta_i}{\gamma_i}(K_i-K_i^*)\dot K_i\notag\\
&= -\beta_i K_i\sigma_i+s_i\tilde d_i
+\frac{\beta_i}{\gamma_i}(K_i-K_i^*)\!
\left[\gamma_i\sigma_i b(K_i)-c_i\!\Big(\tfrac{K_i}{K_{\max,i}}\Big)^{\!3}
-\sigma^{\!*}_i(K_i-0.95K_{i0})g_i\right].
\label{eq:Vdot_raw}
\end{align}
The first two bracketed pieces reproduce the position note exactly. Grouping the
$K_i\sigma_i$ terms via the $\beta_i$-match,
$-\beta_iK_i\sigma_i+\beta_iK_i\sigma_i b(K_i)=-\beta_iK_i\sigma_i(1-b(K_i))$,
\begin{equation}
\dot V_i=
-\beta_i\!\left[K_i\sigma_i(1-b(K_i))+K_i^*\sigma_i b(K_i)\right]
+s_i\tilde d_i
-\frac{\beta_i c_i}{\gamma_i}(K_i-K_i^*)\Big(\tfrac{K_i}{K_{\max,i}}\Big)^{\!3}
\;\underbrace{-\;\frac{\beta_i\sigma^{\!*}_i}{\gamma_i}(K_i-K_i^*)(K_i-0.95K_{i0})g_i}_{\text{leak contribution }L_i}.
\label{eq:Vdot}
\end{equation}

\subsection{The leak term is benign}
$L_i=-\dfrac{\beta_i\sigma^{\!*}_i}{\gamma_i}(K_i-K_i^*)(K_i-0.95K_{i0})g_i$ is a
\emph{product of two affine factors} in $K_i$, with the positive weight
$\beta_i\sigma^{\!*}_i g_i/\gamma_i$. Two cases:
\begin{itemize}[leftmargin=1.5em]
\item If $K_i^*$ and $0.95K_{i0}$ lie on the same side of $K_i$, the two factors share
a sign and $L_i\le0$: the leak \emph{helps} $\dot V_i$.
\item Otherwise $L_i>0$ but is uniformly bounded. The product of affine factors is
maximised at the midpoint $K_i=\tfrac12(K_i^*+0.95K_{i0})$, giving
\begin{equation}
L_i\le\frac{\beta_i\sigma^{\!*}_i\,\bar g_i}{4\gamma_i}\,(K_i^*-0.95K_{i0})^2,
\qquad \bar g_i=e^{\kappa}=e^{0.25}\approx1.284,
\label{eq:Lbound}
\end{equation}
using $g_i\le e^{\kappa}$ (the exponent is largest at $s_i=0$). This is a fixed
additive constant, exactly like the cubic-pushback residual: it enlarges the ultimate
ball but does not threaten sign-definiteness of the dominant $-\beta_iK_i^*\sigma_i$
term.
\end{itemize}
Numerically, with $\sigma^{\!*}_x=0.5$, $\gamma_x=8$, $\beta_x=0.020$, and a generous
gap $|K_x^*-0.95K_{x0}|\le50$,
\[
L_x\le\frac{0.020\cdot0.5\cdot1.284}{4\cdot8}(50)^2\approx1.0\ \text{(m/s}^2\text{ units)},
\]
comparable to the cubic-pushback residual and small against the dominant term over
the operating range. The leak's \emph{equilibrium} (where its own contribution to
$\dot K_i$ vanishes) is $K_i=0.95K_{i0}$, which is the warm-start anchor.

\section{Region 1: adaptation active ($K_i\ll\Kmax-2$)}
\label{sec:region1}
Here $b(K_i)\approx1$, so $1-b(K_i)\approx0$ and \eqref{eq:Vdot} reduces to
\begin{equation}
\dot V_i\le-\beta_iK_i^*\sigma_i+s_i\tilde d_i
+\underbrace{\frac{\beta_ic_i}{\gamma_i}K_i^*\Big(\tfrac{K_i^*}{\Kmax}\Big)^{\!3}}_{\le\,\beta_ic_i\Kmax/\gamma_i}
+\underbrace{\frac{\beta_i\sigma^{\!*}_i\bar g_i}{4\gamma_i}(K_i^*-0.95K_{i0})^2}_{\text{leak residual }\eqref{eq:Lbound}} .
\label{eq:region1}
\end{equation}
Outside the boundary layer ($|s_i|\ge\varepsilon$) the smoothed sign gives
$\sigma_i\ge|s_i|\tanh1\approx0.7616|s_i|$, hence
\begin{equation}
\dot V_i\le-\big(\beta_i\tanh1\,K_i^*-\Deff_i\big)|s_i|+\text{(const)} .
\end{equation}
The first term is strictly negative when $\beta_i\tanh1\,K_i^*>\Deff_i$. Using
$\Deff_i=\beta_i\Dwr_i$ (both $\dot s_i$ and $\tilde d_i$ are projected by the same row
of $\Maug^{-1}$), the $\beta_i$ cancels and the condition is on the \emph{raw} gain:
\begin{equation}
\boxed{\;K_i^*>\frac{\Dwr_i}{\tanh1}\approx1.313\,\Dwr_i\;}
\label{eq:B1raw}
\end{equation}
--- identical in form to the position note. The leak residual and cubic residual
both enter only as fixed additive constants, so they set the size of the ultimate
ball, not its existence. Inside the boundary layer $\Tr{s_i/\varepsilon}\approx
s_i/\varepsilon$ and the surface linearises to
$\dot s_i\approx-\tfrac{\beta_iK_i}{\varepsilon}s_i+\tilde d_i$, attracting $s_i$ to
\begin{equation}
|s_i|_{\mathrm{ss}}\lesssim\frac{\varepsilon\,\Dwr_i}{K_i},
\label{eq:ss}
\end{equation}
with $\beta_i$ again cancelling through $\Deff_i=\beta_i\Dwr_i$.

\section{Region 2: gain saturating ($K_i\to\Kmax$)}
\label{sec:region2}
As $K_i$ passes $\Kmax-2$, $b(K_i)\to0$ and the growth term shuts off. From
\eqref{eq:gainlaw} the gain rolls back via cubic pushback and leak,
$\dot K_i\approx-c_i(K_i/\Kmax)^3-\sigma^{\!*}_i(K_i-0.95K_{i0})g_i<0$, so the cap is
asymptotic, not penetrated. With $b\approx0$ in \eqref{eq:Vdot} the first bracket
becomes $K_i\sigma_i\ge(\Kmax-2)\sigma_i$, and outside the boundary layer
\begin{equation}
\dot V_i\le-\big(\beta_i\tanh1(\Kmax-2)-\Deff_i\big)|s_i|+\text{(const)},
\end{equation}
negative provided $\Dwr_i<\tanh1(\Kmax-2)$. Region 2 is therefore the
\emph{safest} region for stability; its cost is sustained actuator saturation and
chatter, which is exactly what the cubic pushback and leak exist to avoid.

\section{Uniform ultimate boundedness}
\label{sec:uub}
\begin{theorem}[UUB of $(s,K)$ for the velocity controller]
Under the bounded-disturbance assumption $|[\Maug^{-1}d]_i|\le\Deff_i$, choose per axis
\begin{equation}
\Kmax>\frac{\Dwr_i}{\tanh1}\approx1.313\,\Dwr_i,
\qquad
K_i^*\in\Big(\tfrac{\Dwr_i}{\tanh1},\,\Kmax-2\Big),
\end{equation}
and any $K_i(0)\ge0$. Then there is a ball
\[
\mathcal{B}=\Big\{(s,K):|s_i|\le\varepsilon+O(\Dwr_i/\Kmax),\;
0\le K_i\le\Kmax+O\big((c_i+\sigma^{\!*}_i)/\gamma_i\big)\Big\}
\]
that every trajectory of \eqref{eq:canonical},\eqref{eq:gainlaw} enters in finite time
and never leaves.
\end{theorem}
\begin{proof}[Sketch]
$V$ is positive definite, radially unbounded, $C^1$. Outside the boundary layer in
Region 1, summing \eqref{eq:region1} over axes gives
$\dot V\le-\sum_i\beta_i(\tanh1\,K_i^*-\Dwr_i)|s_i|+\sum_i\big(\beta_ic_i\Kmax/\gamma_i
+L_i^{\max}\big)$; the first sum is strictly negative under \eqref{eq:B1raw} and the
second is a fixed positive constant, so $V$ enters and stays in a sublevel set. Inside
the boundary layer $|s_i|\le\varepsilon$ bounds $s$ directly; Region 2 bounds the
$K$-component because both cubic pushback and the leak prevent $K_i$ from exceeding
$\Kmax$ by more than $O((c_i+\sigma^{\!*}_i)/\gamma_i)$. UUB follows by the standard
Lyapunov--Yoshizawa argument.
\end{proof}

\section{Disturbance headroom for the velocity configuration}
\label{sec:headroom}
The admissible task-space disturbance per axis is $\Dwr_i<\tanh1\cdot\Kmax$. With the
velocity caps $\Kmax=(60,80,100)\,$N\,m,
\begin{center}
\begin{tabular}{lccc}
\toprule
Axis $i$ & $x$ & $y$ & $\psi$\\
\midrule
$\Kmax$ (velocity, no-DOB) & 60 & 80 & 100\\
Admissible $\Dwr_i=\tanh1\cdot\Kmax$ (N\,m) & 45.7 & 60.9 & 76.2\\
Cross-coupling cost (Rem.~\ref{rem:coupling}, at $\Kmax$) & 8.1 & 6.7 & 2.8\\
\bottomrule
\end{tabular}
\end{center}

\subsection{The binding constraint moves to actuator saturation}
The dominant high-gain disturbance is the actuator-saturation residual (B5 of the
position note). The per-wheel switching torque after the wheel-map, in the
worst-case aligned configuration with all three gains at the cap and all surfaces
outside their boundary layers, is
\[
|M^{\mathrm{sw,wheel}}|\le\tfrac14\big(\Kmax^x+\Kmax^y+L_\psi\Kmax^\psi\big)\tanh1
\approx0.25(60+80+0.13\cdot100)(0.762)\approx29.1\ \text{N\,m},
\]
with $L_\psi=R/(l+h)\approx0.13$. Against $M_{\max}=10\,$N\,m the per-wheel residual is
$29.1-10\tanh(2.91)\approx19.2\,$N\,m, projecting to
$\Delta D^{\mathrm{sat}}\le4\times19.2\approx76.7\,$N\,m on a single task axis.

\begin{remark}[Saturation overruns $x$-axis headroom at the simultaneous-cap worst case]
This is the one place the velocity configuration is materially \emph{tighter} than the
position case. With the position note's $\Kmax=100$ on all axes the headroom was
$76.2\,$N\,m everywhere; the velocity $x$-axis cap of $60$ drops its headroom to
$45.7\,$N\,m, while the worst-case saturation residual is essentially unchanged at
${\sim}77\,$N\,m. The naive worst-case budget on $x$ therefore reads
\[
\Dwr_x\big|_{\text{worst}}\approx
\underbrace{8.1}_{\text{cross}}+\underbrace{2}_{\text{fric}}+\underbrace{5}_{\text{Coriolis/param}}+\underbrace{76.7}_{\text{sat}}
\approx91.9\ \text{N\,m}\;>\;45.7\ \text{N\,m},
\]
a \emph{negative} margin. The resolution is the same as in the position note: this
combination --- all three $K_i$ pinned at the cap \emph{and} all three surfaces
outside their boundary layers \emph{simultaneously} --- does not persist in operation:
\begin{itemize}[leftmargin=1.5em]
\item The cubic-pushback/leak equilibrium pins $K_i$ at
$K_i^{\mathrm{eq}}=\Kmax(\gamma_i\eta_i/c_i)^{1/3}$ for chatter $\eta_i$, below the cap.
\item Inside the boundary layer $\Tr{s_i/\varepsilon}\approx s_i/\varepsilon\ll1$, so the
switching wrench is far below $K_i$ and saturation never triggers.
\item Saturation is a transient of the reaching phase only.
\end{itemize}
The practical conclusion is sharper than for the position controller: \textbf{with the
velocity caps, the $x$-axis must not be driven outside its boundary layer with $K_x$
near the cap}. This is precisely why $K_{x0}$ should be sized to the \emph{operating}
disturbance (\S\ref{sec:K0}), keeping the reaching phase short and $K_x$ well below the
cap, rather than to the worst case.
\end{remark}

\section{Initial-gain ($K_{0}$) sizing}
\label{sec:K0}
The leak term \eqref{eq:gainlaw} drives $K_i$ toward $0.95K_{i0}$ whenever the
growth term is quiet (small $\sigma_i$). So $K_{i0}$ is the \emph{operating gain the
controller relaxes to}, and it must by itself dominate the \emph{operating} (not
worst-case) disturbance, or the surface escapes its boundary layer before adaptation
can grow $K_i$. The relevant budget excludes the saturation residual (which does not
fire in the boundary layer) and uses the actual small-$K$ cross-coupling
(which scales with the operating $K$, not $\Kmax$). For velocity tracking the
operating disturbance is friction-mismatch $+$ Coriolis/parametric, on the order of
\[
\Dwr_{x,\mathrm{op}}\approx5,\quad
\Dwr_{y,\mathrm{op}}\approx5,\quad
\Dwr_{\psi,\mathrm{op}}\approx2\ \text{N\,m}
\]
(the $\psi$-channel friction spin torque for $\chi=8\,$mm, $\mu=0.6$, $N_i\approx74\,$N,
$V_P\approx0.1\,$m/s is $M_{zi}\approx5.7\times10^{-3}\,$N\,m/wheel, i.e.
${\sim}3\times10^{-3}\,$N\,m on the heading channel --- negligible; the $2\,$N\,m is the
parametric/Coriolis floor). Applying the Region-1 reaching condition \eqref{eq:B1raw}
at startup,
\begin{equation}
\boxed{\;K_{i0}\gtrsim\frac{\Dwr_{i,\mathrm{op}}}{\tanh1}\approx1.313\,\Dwr_{i,\mathrm{op}}\;}
\end{equation}
\begin{center}
\begin{tabular}{lccc}
\toprule
Axis $i$ & $x$ & $y$ & $\psi$\\
\midrule
Operating $\Dwr_{i,\mathrm{op}}$ (N\,m) & 5.0 & 5.0 & 2.0\\
Lower bound $K_{i0}\ge\Dwr_{i,\mathrm{op}}/\tanh1$ (N\,m) & 6.6 & 6.6 & 2.6\\
Current notebook setting $K_{i0}$ & 5 & 5 & 20\\
\midrule
Equilibrium $K_i^{\mathrm{eq}}=\Kmax(\gamma_i\eta/c_i)^{1/3}$, $\eta=0.01$ & 55.7 & 63.5 & 79.4\\
\bottomrule
\end{tabular}
\end{center}

\paragraph{Reading the table.}
\begin{itemize}[leftmargin=1.5em]
\item \textbf{$x,y$:} the current $K_{0}=5$ is marginally \emph{below} the
$6.6\,$N\,m floor. At exactly $5$ the startup surface can drift outside the boundary
layer for the first ${\sim}50\,$ms until adaptation lifts $K$. Raising to
$K_{x0}=K_{y0}=8$--$10$ buys a clean ${\sim}1.5$--$2\times$ margin over the operating
disturbance with no downside (the cubic/leak still cap the steady gain at
$K^{\mathrm{eq}}$).
\item \textbf{$\psi$:} the current $K_{\psi0}=20$ is far \emph{above} the $2.6\,$N\,m
floor --- a deliberate, defensible warm start. Heading needs aggressive initial
authority to suppress the DOB-onset transient (the structural translation/yaw
asymmetry: translation compensation gain ${\sim}2.25\times$ vs.\ yaw lever
$L_\psi\approx0.13$). Sitting $K_{\psi0}$ near $0.2\Kmax$ pre-loads the heading loop so
the reaching phase is essentially instantaneous, at the cost of a little startup
chatter --- exactly the trade the leak term is tuned for. The same asymmetry is the DC
bias source that drives the heading drift of \S\ref{sec:frame}: a high $K_{\psi0}$
shrinks the steady-state $s_\psi^{\mathrm{ss}}=\varepsilon_\psi\Dwr_\psi/K_\psi$ and so
reduces the open-loop drift rate, but it cannot eliminate it --- only the outer heading
loop \eqref{eq:outer} does that. No change recommended, but $K_{\psi0}$ is not a
substitute for closing the cascade.
\end{itemize}

\paragraph{Cascade terms in the operating budget.} The two cascade disturbances of
\S\ref{sec:cascade} --- the setpoint-rate term ($\le k_{p,\psi}\bar s_\psi$, a few
percent) and the heading-rotation term (${\sim}1.35\,$N\,m on translation) --- are
already inside the $\Dwr_{i,\mathrm{op}}$ figures above; neither moves the $K_{i0}$
floors materially.

\paragraph{Reaching time at startup.} From the quadratic balance
$t_{\mathrm{reach}}\lesssim \Dwr_i/(\gamma_i|s_i(0)|)$, with $|s_i(0)|\sim$ the initial
velocity error: for $\Dwr\sim5\,$N\,m, $\gamma_x=8$, $|s_x(0)|\sim0.5\,$m/s, reaching
the boundary layer takes ${\sim}5/(8\cdot0.5)\approx1.25\,$s if starting from
$K=0$ --- but only ${\sim}0.2\,$s when warm-started at $K_{x0}=8$, which already nearly
satisfies \eqref{eq:B1raw}. This is the concrete payoff of sizing $K_{0}$ correctly.

\section{Mixed relative-degree variant: degree-1 translation, degree-2 yaw}
\label{sec:mixed}
The cascade of \S\ref{sec:cascade} bounds heading with an \emph{outer} loop. An
alternative --- preferred for deployment and for generating PINN data that matches the
controller actually used at run time --- folds the heading objective \emph{into the
ASMC itself} by giving the yaw axis its natural relative degree. The translational axes
stay exactly as analysed (degree 1, velocity error); only the yaw surface changes:
\begin{equation}
s_x=\Vx-\Vx^d,\quad s_y=\Vy-\Vy^d,\qquad
\boxed{\,s_\psi=\dot e_\psi+\lambda_\psi(e_\psi)\,e_\psi\,},
\label{eq:mixed_surface}
\end{equation}
with $e_\psi$ the shaped heading error \eqref{eq:epsi} and $\dot e_\psi=e_\psi'(d_\psi)
\,(\dot\psi-\dot\psi_d)$. This is the \emph{exact} yaw surface of \texttt{asmc\_torques};
the mixed controller is ``velocity ASMC on $(x,y)$, position ASMC on $\psi$.'' The
relative-degree structure now matches the control objective axis-by-axis: translation
wants velocity (degree 1, position free), heading wants the angle (degree 2). The
asymmetry is correct, not a hack.

\subsection{Why this kills the drift without an integrator}
The degree-1 yaw residual integrated into drift, $\Delta\psi=\int s_\psi\,d\tau$
\eqref{eq:dpsi_integral}. With \eqref{eq:mixed_surface}, the manifold $s_\psi=0$ is
itself the stable first-order law $\dot e_\psi=-\lambda_\psi e_\psi$, so $e_\psi\to0$
\emph{on the surface}: a biased $\tilde d_\psi$ no longer integrates, it is absorbed
into a bounded steady error exactly as in the position note,
\begin{equation}
|e_\psi|_{\mathrm{ss}}\lesssim\frac{s_\psi^{\mathrm{ss}}}{\lambda_{\psi,\min}}
=\frac{\varepsilon_\psi\,\Dwr_\psi}{\lambda_{\psi,\min}\,K_\psi}.
\label{eq:mixed_ess}
\end{equation}
\begin{center}
\begin{tabular}{lcc}
\toprule
& $\Dwr_\psi=1.0$ & $\Dwr_\psi=2.0$\\
\midrule
$|e_\psi|_{\mathrm{ss}}$ at $\lambda_{\psi,\min}=0.5$ (deg) & 0.116 & 0.231\\
$|e_\psi|_{\mathrm{ss}}$ at $\lambda_{\psi,\max}=5.0$ (deg) & 0.012 & 0.023\\
\bottomrule
\end{tabular}
\end{center}
The dynamic Gaussian $\lambda_\psi(e_\psi)$ rises toward $\lambda_{\psi,\max}$ as
$e_\psi\to0$, so the \emph{operating} bound is the tight right column
($\le0.02^\circ$); the worst case uses $\lambda_{\psi,\min}$. Either way it is a fixed
offset (Gap 1, benign), versus the degree-1 secular drift of $1.7$--$3.5^\circ$ over
$30\,$s. No $k_{p,\psi}$, no timescale-separation assumption --- $\lambda_\psi$ does the
job a single self-contained certificate.

\subsection{What the proof inherits, and what changes}
The aggregate Lyapunov function becomes \emph{mixed}: the $(x,y)$ blocks keep the
relative-degree-1 canonical form of \S\S\ref{sec:plant}--\ref{sec:uub} verbatim, while
the $\psi$ block reverts to the position note's degree-2 analysis --- the dynamic-$\lambda$
cancellation $\dot s_\psi=[\Maug^{-1}\tau_{sw}]_\psi+[\Maug^{-1}d]_\psi$ holds because
the yaw equivalent control again carries the $\dot\lambda_\psi$ feedforward (below). The
per-axis $\beta_i$-cancellation and the (B1) bound $\Kmax>\Dwr_\psi/\tanh1$ are
unchanged; one simply can no longer claim ``uniformly relative degree 1.'' Two genuine
shifts:
\begin{itemize}[leftmargin=1.5em]
\item \emph{Reduced-order heading dynamics.} On the manifold the heading settles with
time constant $1/\lambda_\psi\in[0.2,2.0]\,$s. The reaching (boundary-layer) bandwidth
$\beta_\psi K_\psi/\varepsilon_\psi\approx171\,$rad/s is unchanged.
\item \emph{Cross-coupling carries a position-type term.} The yaw switching
$K_\psi\tanh(s_\psi/\varepsilon_\psi)$ still leaks into $x$ via $[\Maug^{-1}]_{x\psi}$,
bounded by $|[\Maug^{-1}]_{x\psi}|K_\psi\approx0.30\,$N\,m ($0.7\%$ of the $x$ headroom)
because $|\tanh|\le1$ --- the Device-1 bound is identical. But $s_\psi$ now contains
$\lambda_\psi e_\psi$ (a position quantity), so during large heading transients the
leak sits near its bound more of the time; the \emph{operating} $\Dwr_x$ during turns is
marginally higher, still negligible against the $\sim$46\,N\,m $x$ headroom.
\end{itemize}

\subsection{The yaw equivalent control must be restored in full}
\label{sec:mixed_eq}
Degree-2 yaw tracking requires the $\dot\lambda_\psi$ feedforward, identical to
\texttt{asmc\_torques}:
\begin{equation}
\alpha_{\mathrm{eq}}=\alpha_{\mathrm{ff}}
-\dot\lambda_\psi\,e_\psi-\lambda_\psi\,\dot e_\psi\,e_\psi',
\label{eq:mixed_alpha}
\end{equation}
where $(\lambda_\psi,\dot\lambda_\psi)=\texttt{get\_dynamic\_lambda}(e_\psi,
e_\psi'\dot e_\psi,\lambda_{\psi,\min},\lambda_{\psi,\max},\mu_\psi)$ and
$\dot e_\psi=e_\psi'(\dot\psi-\dot\psi_d)$. \textbf{One correction that matters
independently of the surface choice:} the present \texttt{asmc\_torques\_vel} yaw
equivalent wrench omits the yaw viscous-drag term that \texttt{asmc\_torques} carries,
\begin{equation}
\Big(\tfrac{4p_1(l+h)^2}{R^2}+\tfrac{8p_2h^2}{(R-R_d)^2}\Big)\dot\psi
\approx 38.2\,\dot\psi\ \text{N\,m},
\label{eq:yaw_visc}
\end{equation}
i.e.\ ${\sim}11.5\,$N\,m at $\dot\psi=0.3\,$rad/s. Left uncancelled this is a large,
$\dot\psi$-proportional component of $\tilde d_\psi$ --- it would \emph{worsen} the very
drift the mixed surface fixes. Restoring it (it is a known, exactly-cancellable drag)
shrinks $\Dwr_\psi$ materially and should be done whether or not the surface is changed.

\subsection{Code edits for the mixed controller}
\label{sec:mixed_code}
Relative to the current \texttt{asmc\_torques\_vel}, the edits are confined to the yaw
channel; $(x,y)$ are untouched.
\begin{enumerate}[leftmargin=2em]
\item \textbf{Unpack pose and read the desired heading.} Add $\psi=\texttt{state[4]}$;
compute $d_\psi=\psi-\texttt{get\_psi\_des(t)}$, then $e_\psi$ and $e_\psi'$ from
\eqref{eq:epsi} (the \texttt{e\_psi}/\texttt{edash\_psi} expressions).
\item \textbf{Degree-2 yaw surface.} Replace $s_\psi=\dot\psi-\omega_d$ with
$\dot e_\psi=e_\psi'\,(\dot\psi-\texttt{get\_omega\_des\_v(t)})$ and
$s_\psi=\dot e_\psi+\lambda_\psi\,e_\psi$, where $(\lambda_\psi,\dot\lambda_\psi)=
\texttt{get\_dynamic\_lambda}(e_\psi,e_\psi'\!\cdot\!(\dot\psi-\omega_{\mathrm{ff}}),
\dots)$. Leave $s_x,s_y$ exactly as they are.
\item \textbf{Yaw equivalent control \eqref{eq:mixed_alpha}.} Replace
$\alpha_{\mathrm{eq}}=\texttt{get\_alpha\_des\_v(t)}$ with
$\alpha_{\mathrm{eq}}=\alpha_{\mathrm{ff}}-\dot\lambda_\psi e_\psi-\lambda_\psi\dot
e_\psi e_\psi'$.
\item \textbf{Restore the yaw viscous term \eqref{eq:yaw_visc}} by adding
$+\big(4p_1(l+h)^2/R^2+8p_2h^2/(R-R_d)^2\big)\dot\psi$ inside \texttt{M\_psi\_eq}, matching
\texttt{asmc\_torques}.
\item \textbf{Parameters.} The yaw $\lambda$ fields already exist on
\texttt{ASMCParams} ($\lambda_{\psi,\min}{=}0.5$, $\lambda_{\psi,\max}{=}5.0$,
$\mu_\psi{=}100$); no \texttt{kp\_psi} is needed (the cascade gain is retired). The
adaptive-gain block and the $(x,y)$ surfaces are unchanged.
\end{enumerate}
At $t=0$, $\psi(0)=\psi_d(0)\Rightarrow e_\psi(0)=0$ and $\lambda_\psi e_\psi(0)=0$, so
the yaw surface starts on its manifold with no startup transient, provided
$\dot\psi(0)=\omega_{\mathrm{ff}}(0)$ (already true by construction).

\subsection{PINN-data implication}
The mixed controller changes the yaw excitation the network sees: $s_\psi$ carries
$\lambda_\psi e_\psi$, so the heading-channel torque spectrum on turns gains
low-frequency content (the $\lambda_\psi e_\psi$ term) absent from the degree-1 data.
This is the controller that will run under MPC, so it is the correct distribution to
train on --- but datasets generated with the degree-1 velocity controller and the
mixed controller must \emph{not} be pooled; regenerate consistently.

\section{Conclusion}
The central result of this revision is a correction: body-frame velocity-error UUB does
\emph{not} certify the intended trajectory. A bounded pose offset is harmless for PINN
training (Gap 1), but a \emph{growing} heading error (Gap 2) changes the trajectory
class --- a circle opens into a spiral. Because heading error is the time-integral of
the yaw surface, $\Delta\psi=\int s_\psi\,d\tau$, any DC bias in the yaw disturbance
(the $\chi^2$ spin torque, phase-locked roller handoff, and above all the
$2.25{:}0.13$ DOB translation/yaw asymmetry) integrates into linear drift ---
quantified here at $2$--$5^\circ$ per circle lap. The inner-loop UUB proof is silent on
this. The fix is the cascade: an outer heading loop $\omega^d=\dot\psi_d+k_{p,\psi}
e_\psi(d_\psi)$, acting on the position controller's \emph{shaped} heading error
$e_\psi$, converts the unbounded integrator into a stable nonlinear first-order system,
bounding $d_\psi$ at a fixed sub-$0.1^\circ$ offset. Because $e_\psi'\ge1$ everywhere,
the shaped form is never slower than a raw-error correction and recovers faster from
large excursions. Timescale separation is
comfortable (inner $25$--$171\,$rad/s; outer $\le5$--$34\,$rad/s), and the cascade's
added disturbance is negligible against the (B1) headroom.

For deployment and for PINN data that matches the run-time controller, the
\emph{mixed relative-degree} variant (\S\ref{sec:mixed}) is preferred: keep the
degree-1 velocity surfaces on $(x,y)$ but give yaw its natural degree-2 surface
$s_\psi=\dot e_\psi+\lambda_\psi e_\psi$ on the shaped error. The manifold itself then
enforces $\dot e_\psi=-\lambda_\psi e_\psi$, so a biased $\tilde d_\psi$ produces a
\emph{bounded} $|e_\psi|_{\mathrm{ss}}\le\varepsilon_\psi\Dwr_\psi/(\lambda_{\psi,\min}
K_\psi)\lesssim0.02$--$0.23^\circ$ rather than secular drift --- no outer loop, no
timescale-separation caveat. The cost is a mixed certificate (degree-1 on $(x,y)$,
degree-2 on $\psi$) and the restored yaw equivalent control, including the
$\sim38.2\,\dot\psi\,$N\,m viscous term \eqref{eq:yaw_visc} the present velocity
controller omits --- a correction worth making on its own merits, since uncancelled it
is itself an $\sim$11\,N\,m yaw disturbance.

The relative-degree-one inner-loop analysis itself transfers cleanly from the position
note: no dynamic $\lambda$, no $\dot\lambda$ feedforward, the canonical
$\dot s=-\tilde K\tanh(s/\varepsilon)+\tilde d$ falls out directly, and the
$\beta_i$-cancellation and (B1) $\Kmax>\Dwr/\tanh1$ hold unchanged. Three numerical
facts:
\begin{enumerate}[leftmargin=2em]
\item Inverting the augmented $\Maug$ lowers worst-axis cross-coupling to
${\sim}13.4\%$ (from ${\sim}17.7\%$).
\item The smaller velocity caps make the simultaneous-cap actuator-saturation residual
overrun the $x$-axis headroom; the design response is to keep $K_x$ off the cap by
sizing $K_{x0}$ to the operating disturbance, not the worst case.
\item The leak-to-$K_0$ term is sign-indefinite but benign --- a product of two affine
factors bounded by \eqref{eq:Lbound} --- pinning the relaxed gain at $0.95K_{i0}$.
\end{enumerate}
The initial-gain rule $K_{i0}\gtrsim\Dwr_{i,\mathrm{op}}/\tanh1$ recommends nudging
$K_{x0},K_{y0}$ from $5$ to ${\sim}8$--$10\,$N\,m and confirms $K_{\psi0}=20$ as an
intentional heading warm start --- though a high $K_{\psi0}$ only \emph{reduces} the
drift rate; closing the heading loop is what \emph{bounds} it.

\appendix
\section{Numerical scratchpad}
\label{sec:scratch}
$\tanh1=(e^2-1)/(e^2+1)=0.76159\ldots$\quad
$L_\psi=R/(l+h)=0.05/0.385=0.12987$.\quad
$\tilde m=m_s+4J_1/R^2=35.6+4(5.87\times10^{-3})/0.0025=44.992$.\quad
$I_{\psi,\mathrm{aug}}=I_s+4(l+h)^2J_1/R^2=4.42+4(0.385)^2(5.87\times10^{-3})/0.0025=5.812$.\quad
Smallest eigenvalue of $\Maug\approx5.79$.\quad
$\beta_x\le[\Maug^{-1}]_{xx}=0.02228$, $\beta_y\le0.02225$,
$\beta_\psi\le0.17261$.\quad
Leak modulation cap $\bar g=e^{0.25}=1.2840$; leak vanishes at $K=0.95K_0=(4.75,4.75,19.0)$.

\end{document}
